
\chapter{Пользовательский интерфейс}
\label{chap:user-interface}

CIAO позволяет создавать приложения с графическим
интерфейсом на основе Tkinter.

При этом логика автоматов отделяется от элементов
интерфейса. Кнопки отправляют команды, а автоматы
передают данные, которые необходимо отобразить.

Рассмотрим организацию интерфейса на примере
Array Sum.

\section{Взаимодействие с интерфейсом}

В приложении используются три основных компонента:

\begin{itemize}
  \item \code{Controller} выполняет вычисления
    и управляет их последовательностью.
  \item \code{ViewAutomaton} связывает вычисления
    с пользовательским интерфейсом.
  \item \code{Window} содержит Tkinter-виджеты
    и отвечает за их отображение.
\end{itemize}

View объявляется наследником \code{ViewAutomaton}
и работает с командами так же, как обычный автомат.

\section{Привязка Window к View}

Для взаимодействия с автоматом Window предоставляет
кнопки, методы чтения пользовательского ввода и
методы отображения данных.

View, в свою очередь, объявляет два вида внешних операций:
\code{provided} для приёма событий и \code{required}
для передачи данных интерфейсу.

Их связывает \code{TkConnection}.

\begin{figure}[ht]
\centering
\begin{tikzpicture}[
  component/.style={
    draw=ciaoRule,
    fill=ciaoCodeBg,
    rounded corners=1mm,
    minimum width=40mm,
    minimum height=22mm,
    align=center,
    font=\small\sffamily
  },
  flow/.style={
    -{Stealth[length=2mm]},
    draw=ciaoBlue,
    thick
  }
]
  \node[component] (window) at (0,0) {
    \textbf{Window}\\
    init\_button\\
    render\_data()
  };

  \node[component] (view) at (10,0) {
    \textbf{SumView}\\
    init\_clicked\\
    render\_data
  };

  \draw[flow]
    ([yshift=5mm]window.east) --
    node[above,font=\scriptsize] {bind\_event}
    ([yshift=5mm]view.west);

  \draw[flow]
    ([yshift=-5mm]view.west) --
    node[below,font=\scriptsize] {bind\_command}
    ([yshift=-5mm]window.east);
\end{tikzpicture}
\caption{Два направления привязки Window и View}
\label{fig:ui-binding}
\end{figure}

\subsection{От Window к View}

Привяжем кнопку запуска к внешней provided-команде
\code{SumView.init_clicked}:

\begin{lstlisting}[language=Python,numbers=none]
connection.bind_event(
    window.init_button,
    (sample.view, SumView.init_clicked),
    payload_source=FactoryPayloadSource(
        StartInput, window.read_start_input
    ),
)
\end{lstlisting}

При нажатии кнопки \code{FactoryPayloadSource}
вызывает \code{window.read_start_input()} и формирует
объект \code{StartInput}.

Затем \code{TkConnection} отправляет команду
\code{init_clicked} автомату View через Runtime.
View обрабатывает её как обычное событие автомата
и при необходимости передаёт команды Controller.

\subsection{От View к Window}

Для обновления интерфейса View объявляет внешнюю
required-команду \code{render_data}. Она связывается
с методом Window:

\begin{lstlisting}[language=Python,numbers=none]
connection.bind_command(
    (sample.view, SumView.render_data),
    window.render_data,
)
\end{lstlisting}

Когда View отправляет \code{render_data} через Emit,
\code{TkConnection} вызывает зарегистрированный
метод \code{window.render_data(...)} с переданными
данными.

Window обновляет соответствующие виджеты,
не обращаясь к контексту автомата.

Таким же образом можно подключить команды
\code{render_controls} и \code{render_status}
для управления кнопками и отображения статуса.

\begin{importantbox}
Обработчики кнопок не назначаются напрямую через
Tkinter \code{command=}. Их регистрирует
\code{TkConnection}.

На одну кнопку допускается одна CIAO-привязка,
а на внешнюю required-операцию --- не более
одного обработчика.
\end{importantbox}


\section{Исполнение и обработка ошибок}

При нажатии кнопки \code{TkConnection} отправляет
команду в Runtime и планирует её обработку через
механизм \code{after()} библиотеки Tkinter.

Каждый запланированный вызов выполняет один
\code{Runtime.step()}. Это позволяет обрабатывать
очередь постепенно, сохраняя возможность
обновлять интерфейс между шагами.

Если во время исполнения возникает ошибка,
\code{TkConnection} останавливает обработку
и передаёт сведения о ней в
\code{BaseWindow.show_error}.

По умолчанию окно может показать сообщение об ошибке,
но этот метод можно переопределить для собственного
способа отображения.

\begin{notebox}
Tkinter-виджеты необходимо обновлять из потока,
в котором работает окно. Не запускайте
\code{Runtime.run()} непосредственно из обработчиков
кнопок, подключённых через \code{TkConnection}.

Если подключение завершается до закрытия окна,
освободите его с помощью \code{dispose()}.
\end{notebox}
